% language=us runpath=texruns:manuals/workflows

\environment workflows-style

\startcomponent workflows-embedding

\startchapter[title={Embedding text}]

When you embed graphics that use fonts the filesize can be reduced by sharing
fonts. Doing this is somewhat complicated by the fact that the logic must be
robust which is not trivial given the way \PDF\ files can be made. For instance,
when I implemented this feature in \MKIV\ I had an application in mind that in
the end was not suitable: all the included text in the (\EPS\ images) turned out
to be either unfamiliar \TYPEONE\ clones (operating systems and graphic editors
can make mess of this) or the text was turned into vectors images (probably to
work around possible validation issues).

So, for now I assume that we're dealing with for instance text made by \TEX\
(read: \LUATEX), for instance typesetting buffers in a sub|-|run (which is handy
when writing a manual) or using precooked \METAFUN\ images. In order for this to
work (at least for now) we need to enable inclusion via \LUA: \footnote {In
\CONTEXT\ \LMTX\ these things will happen automatically.}

\startbuffer
\enabledirectives[graphics.pdf.uselua]
\stopbuffer

\typebuffer \getbuffer

Also, because we don't want unexpected side effects, you need to be explicit
about sharing fonts using the \type {compact} option, as in:

\starttyping
\externalfigure[one.pdf][compact=yes]
\externalfigure[two.pdf][compact=yes]
\stoptyping

Here we give an example of embedded typeset buffers:

\startbuffer
\startbuffer[one]
\setupbodyfont[dejavu]
\startTEXpage
  \samplefile{ward} \sl \samplefile{klein} \bf \samplefile{sapolsky}
\stopTEXpage
\stopbuffer

\startbuffer[two]
\setupbodyfont[dejavu]
\startTEXpage
  \samplefile{klein} \bf \samplefile{ward} \sl \samplefile{sapolsky}
\stopTEXpage
\stopbuffer
\stopbuffer

\typebuffer \getbuffer

The actual typesetting takes place with:

\startbuffer
\startcombination
  {\typesetbuffer[one]
    [height=.25\textheight,compact=yes,frame=on,rulethickness=.2pt]}
  {one}
  {\typesetbuffer[two]
    [height=.25\textheight,compact=yes,frame=on,rulethickness=.2pt]}
  {two}
\stopcombination
\stopbuffer

\typebuffer

In \in {figure} [fig:sharedfonts-1] we see the two pages. The \type {compact}
option makes sure that we share fonts when possible. In this case we used
\LUATEX\ to create these (independent) \PDF\ files so we can be rather sure that
we use the same fonts. On the console the merging is reported:

\starttyping
mkiv lua stats  > merged fonts:
    DejaVuSerif (+38)
    DejaVuSerif-Bold (+39)
    DejaVuSerif-Italic (+39)
\stoptyping

The number of characters that get forced into the file is shown between
parentheses but the real number can be larger of course when the font is also
used in the main document.

\startplacefigure[title={A few typeset pages.},reference=fig:sharedfonts-1]
  \getbuffer
\stopplacefigure

A second example involves \METAPOST\ graphics. Here we use buffers but you can
imagine collecting graphics in one file and then including them with:

\starttyping
\externalfigure[myimages.pdf][compact=yes,page=4]
\stoptyping

Beware: this only makes sense if such graphics take runtime, which in a normal
usage of \METAPOST\ in \CONTEXT\ is no real issue: runtime rendering is quite
efficient and not slower than inclusion. The example code is:

\startbuffer
\startbuffer[three]
\setupbodyfont[dejavu]
\starttext
\startMPpage
  draw textext("Example 1") xsized 3cm withcolor "darkred" ;
  draw boundingbox currentpicture enlarged 2mm withcolor "darkblue" ;
\stopMPpage
\startMPpage
  draw textext("Example 2") xsized 3cm withcolor "darkblue" ;
  draw boundingbox currentpicture enlarged 2mm withcolor "darkred" ;
\stopMPpage
\stoptext
\stopbuffer
\stopbuffer

\typebuffer \getbuffer

\startbuffer
\startcombination
  {\typesetbuffer[three][compact=yes,page=1]} {Page 1}
  {\typesetbuffer[three][compact=yes,page=2]} {Page 2}
\stopcombination
\stopbuffer

\typebuffer

\startplacefigure[title={A few prepared graphics.},reference=fig:sharedfonts-2]
  \getbuffer
\stopplacefigure

In \in {figure} [fig:sharedfonts-2] we see the two graphics. In this case they
come from one file so the fonts are shared in embedding anyway, but the \type
{compact} option will make sure that they are also shared with the main document
fonts if possible. In fact, while normal embedding will include the normal font
resources the \type {compact} variant will only embed the glyphs that are used,
so it might even be more efficient. Because the examples use digits we actually
get this report on the console now:

\starttyping
mkiv lua stats  > merged fonts:
    DejaVuSerif (+40)
    DejaVuSerif-Bold (+39)
    DejaVuSerif-Italic (+39)
\stoptyping

This mechanism might be refined, depending on realistic use cases, user input and
my own demands. Currently we use the font name as check but we might add some kind
of mapping mechanism to work around fuzzy naming of what is actually the same
font.

There is more possible when embedding files, but that is discussed in the \type
{pdfmerge} manual.

\stopchapter

\stopcomponent
